\bmtchapitre{Étude complète et représentation graphique}{Assembler domaines, limites, dérivée, points remarquables et positions relatives pour tracer une courbe justifiée.}{Analyse}
\label{t11s-c05}
\begin{revision}Les quatre chapitres précédents; signe d’un polynôme; équations de droites.\end{revision}
\begin{automatismes}\exo\label{t11s-c05-auto}Factoriser $x^2-2x-3$. Calculer $f(0)$ pour $f(x)=x+1+2/(x-1)$.\end{automatismes}
\begin{corrige}[exercice~\ref{t11s-c05-auto}]\label{sol-t11s-c05-auto}$(x-3)(x+1)$; $f(0)=-1$.\end{corrige}

\section{Organiser l’étude}
\begin{methode}{une étude en six étapes}
Domaine et éventuelle réduction par symétrie/périodicité; limites et branches infinies; dérivée et variations; points d’intersection avec les axes; tangentes ou points particuliers; tracé cohérent avec l’ensemble des résultats. Pour comparer deux courbes, étudier le signe de leur différence sur leur domaine commun.
\end{methode}
\begin{exemple}[une fraction rationnelle]
Soit $f(x)=\dfrac{x^2+1}{x-1}=x+1+\dfrac2{x-1}$.
Le domaine est $\R\setminus\{1\}$. Les limites en $1^-$ et $1^+$ sont $-\infty$ et $+\infty$. Aux deux infinis, l’asymptote est $y=x+1$.
\[
f'(x)=1-\frac2{(x-1)^2}
=\frac{(x-1)^2-2}{(x-1)^2}.
\]
Elle s’annule en $a=1-\sqrt2$ et $b=1+\sqrt2$. Elle est positive en dehors de $[a;b]$ et négative à l’intérieur, en séparant les deux intervalles du domaine. Les valeurs sont $f(a)=2-2\sqrt2$ et $f(b)=2+2\sqrt2$.
\begin{center}\begin{tabular}{c|ccccccc}
$x$&$-\infty$&&$a$&&$1$&&\\
$f'$&&$+$&$0$&$-$&non défini&&\\
$f$&$-\infty$&$\nearrow$&$2-2\sqrt2$&$\searrow$&$-\infty$&&\\\hline
$x$&$1$&&$b$&&$+\infty$&&\\
$f'$&non défini&$-$&$0$&$+$&&&\\
$f$&$+\infty$&$\searrow$&$2+2\sqrt2$&$\nearrow$&$+\infty$&&
\end{tabular}\end{center}
$f(0)=-1$. Le numérateur $x^2+1$ ne s’annule jamais, donc la courbe ne coupe pas l’axe des abscisses. Pour tout $h\ne0$, $f(1+h)+f(1-h)=4$, et le domaine est symétrique par rapport à $1$ : $(1;2)$ est centre de symétrie.
\end{exemple}
\begin{visualisation}[deux branches, quatre variations]
\begin{center}\begin{tikzpicture}\begin{axis}[width=10cm,height=7cm,axis lines=middle,xmin=-5,xmax=7,ymin=-6,ymax=10,xlabel={$x$},ylabel={$y$},samples=160]
\addplot[bmtSarcelle,thick,domain=-5:0.73]{x+1+2/(x-1)};
\addplot[bmtSarcelle,thick,domain=1.27:7]{x+1+2/(x-1)};
\addplot[bmtOcre,dashed,domain=-5:7]{x+1};
\draw[bmtOcre,dashed](axis cs:1,-6)--(axis cs:1,10);
\end{axis}\end{tikzpicture}\end{center}
Le maximum de la branche gauche et le minimum de la branche droite ne sont pas les mêmes points. L’asymptote verticale sépare le domaine.
\end{visualisation}
\section{Point d’inflexion et tangente}
\begin{definition}[point d’inflexion]Un point d’inflexion est un point où la courbe change de convexité. Pour les fonctions deux fois dérivables étudiées ici, un changement de signe de $f''$ de part et d’autre fournit un critère suffisant (admis). La convexité désigne l’allure courbée vers le haut, la concavité vers le bas.\end{definition}
\begin{exemple}Pour $g(x)=x^3$, $g''(x)=6x$ change de signe en $0$. L’origine est un point d’inflexion, et la tangente $y=0$ est traversée par la courbe. Pour $x^4$, $g''(0)=0$ mais $g''$ reste positive : pas d’inflexion en $0$.\end{exemple}
\begin{avertissement}Annuler $f''$ ne suffit pas. Une courbe apparemment proche d’une droite ne prouve pas une asymptote. Une esquisse est un résumé de résultats déjà justifiés.\end{avertissement}
\section{Comparer deux fonctions}
Si $f-g>0$ sur un intervalle, $\Cf_f$ est au-dessus de $\Cf_g$; si $f-g<0$, elle est en dessous; les zéros de la différence donnent les intersections. Les unités des axes peuvent être différentes; préciser le repère pour les arguments métriques.

% ENRICHISSEMENT C05

\subsection*{Lire une étude comme un ensemble de contraintes}
\begin{methode}{construire la courbe dans le bon ordre}
Tracer d’abord les axes et les asymptotes. Placer ensuite les points connus, en particulier extrema et intersections. Dessiner séparément chaque intervalle du domaine, en respectant les sens de variation et la position par rapport aux droites. Revenir au tableau pour contrôler chaque branche : la calculatrice peut aider à visualiser, mais ne remplace aucune justification.
\end{methode}
\begin{exemple}[compter des intersections sans formule des racines]
Pour $p(x)=x^3-3x$, la dérivée est $3(x-1)(x+1)$. Les intervalles de monotonie stricte sont $]-\infty;-1]$, $[-1;1]$ et $[1;+\infty[$, avec valeurs remarquables $2$ et $-2$. Une horizontale $y=k$ avec $-2<k<2$ rencontre chaque branche exactement une fois : la continuité donne l’existence et la monotonie l’unicité sur chaque branche. Ainsi on peut compter trois solutions sans les écrire en formule explicite.
\end{exemple}
\begin{methode}{confronter une équation au domaine}
Pour $h(x)=(2x+1)/(x-1)$, résoudre $h(x)=c$ exige $x\ne1$. Multiplier alors par $x-1$ donne $(2-c)x=-1-c$. Si $c=2$, on obtient $0=-3$, impossible. Sinon, $x=(1+c)/(c-2)$; ce nombre n’est jamais $1$, car cela imposerait $1=-2$. Le niveau de l’asymptote horizontale n’est pas atteint dans cet exemple, mais cette propriété ne vaut pas pour toute asymptote.
\end{methode}
\begin{avertissement}[les points remarquables n’ont pas le même rôle]
Un extremum local concerne les valeurs voisines; un extremum global concerne tout le domaine. Un point d’inflexion marque un changement de convexité et peut avoir une tangente non horizontale. Une valeur interdite ne devient pas un point de la courbe même si sa limite est finie.
\end{avertissement}

\section{Exercices et problèmes}
\exo[Une parabole]\label{t11s-c05-e01}
Étudier complètement $f(x)=x^2-4x+3$ : domaine, limites, variations, axe et intersections avec les axes. Tracer.
\begin{corrige}[exercice~\ref{t11s-c05-e01}]\label{sol-t11s-c05-e01}
Domaine $\R$, limites $+\infty$ aux deux infinis. $f'=2x-4$; décroissance jusqu’à $2$, croissance ensuite; minimum $-1$. Forme $(x-2)^2-1$, axe $x=2$. Zéros $1,3$; ordonnée à l’origine $3$.
\end{corrige}
\exo[Courbe et droite]\label{t11s-c05-e02}
Comparer $f(x)=x^2-4x+3$ et $g(x)=x-1$.
\begin{corrige}[exercice~\ref{t11s-c05-e02}]\label{sol-t11s-c05-e02}
$f-g=x^2-5x+4=(x-1)(x-4)$. Intersections $(1;0)$ et $(4;3)$. Parabole au-dessus à l’extérieur de $[1;4]$, en dessous sur $]1;4[$.
\end{corrige}
\exo[Un quotient simple]\label{t11s-c05-e03}
Étudier $h(x)=\dfrac{2x+1}{x-1}$ : limites, variations, asymptotes, intersections et centre de symétrie.
\begin{corrige}[exercice~\ref{t11s-c05-e03}]\label{sol-t11s-c05-e03}
$h=2+3/(x-1)$, domaine $x\ne1$; limites $2$ aux infinis, $-\infty$ et $+\infty$ en $1^-,1^+$. $h'=-3/(x-1)^2<0$ sur chaque intervalle. Asymptotes $x=1,y=2$. Intersections $(-1/2;0)$ et $(0;-1)$; centre $(1;2)$ par la somme $h(1+t)+h(1-t)=4$.
\end{corrige}
\exo[Cubicité]\label{t11s-c05-e04}
Étudier $p(x)=x^3-3x$, sa symétrie, ses extrema et son point d’inflexion. Donner sa tangente en $0$.
\begin{corrige}[exercice~\ref{t11s-c05-e04}]\label{sol-t11s-c05-e04}
Domaine $\R$, impaire. $p'=3(x-1)(x+1)$ : maximum local $(-1;2)$, minimum local $(1;-2)$. Limites $-\infty,+\infty$. $p''=6x$ change de signe en $0$; tangente $y=-3x$. Zéros $-\sqrt3,0,\sqrt3$.
\end{corrige}
\exo[Lecture argumentée]\label{t11s-c05-e05}
Un élève affirme : « $f'(a)=0$, donc $f(a)$ est un maximum ». Donner deux contre-exemples différents.
\begin{corrige}[exercice~\ref{t11s-c05-e05}]\label{sol-t11s-c05-e05}
$x^2$ en $0$ a un minimum et dérivée nulle; $x^3$ en $0$ n’a aucun extremum et dérivée nulle. Le signe de part et d’autre est nécessaire pour le critère utilisé.
\end{corrige}
\exo[Oscillation]\label{t11s-c05-e06}
Étudier $u(x)=2\cos x$ sur $[0;\pi]$, puis expliquer comment obtenir sa courbe sur $\R$.
\begin{corrige}[exercice~\ref{t11s-c05-e06}]\label{sol-t11s-c05-e06}
$u'=-2\sin x\le0$ et strictement négative sur $]0;\pi[$. Décroît de $2$ à $-2$; zéro $\pi/2$. La parité complète $[-\pi;\pi]$ et la période $2\pi$ permet les translations.
\end{corrige}

\exo[Le nombre de solutions dépend d’un niveau]\label{t11s-c05-p01}
On reprend $p(x)=x^3-3x$ sur $\R$. À partir de son tableau de variation, déterminer le nombre de solutions distinctes de $p(x)=k$ selon le réel $k$. Traiter séparément $k=-2$ et $k=2$; pour ces deux niveaux, factoriser l’équation et donner les solutions exactes.
\begin{corrige}[exercice~\ref{t11s-c05-p01}]\label{sol-t11s-c05-p01}
Pour $k<-2$ ou $k>2$, il y a une solution; pour $-2<k<2$, il y en a trois. En $k=2$, $x^3-3x-2=(x-2)(x+1)^2$ : deux solutions distinctes $-1,2$. En $k=-2$, $x^3-3x+2=(x+2)(x-1)^2$ : deux solutions distinctes $-2,1$. Les extrema appartiennent à deux intervalles adjacents; il ne faut pas les compter deux fois comme points distincts. Le comptage se justifie par continuité et monotonie sur les branches, pas seulement par le degré trois.
\end{corrige}
\exo[Auditer une description de courbe]\label{t11s-c05-p02}
Pour $f(x)=(x^2+1)/(x-1)$, un compte rendu annonce : « la courbe coupe l’axe des abscisses; ses deux branches se rejoignent en $x=1$; elle possède un maximum global en $1-\sqrt2$ ». Réfuter chacune des trois affirmations à l’aide de l’expression ou des résultats du cours et écrire une description corrigée.
\begin{corrige}[exercice~\ref{t11s-c05-p02}]\label{sol-t11s-c05-p02}
Le numérateur $x^2+1$ est strictement positif, donc aucun zéro. La valeur $1$ est interdite et les limites latérales sont infinies; les branches ne se rejoignent pas. Le point d’abscisse $1-\sqrt2$ est un maximum local de la branche gauche, mais la branche droite tend vers $+\infty$ en $1^+$, donc pas de maximum global. Description : deux branches séparées par $x=1$, aucune intersection avec l’axe horizontal, maximum local à gauche et minimum local à droite, asymptote oblique $y=x+1$.
\end{corrige}

\begin{jemeteste}Pour vérifier mon tracé, je confronte chaque branche aux limites, chaque extremum au tableau et chaque intersection à son équation.\end{jemeteste}
\begin{notepourlaclasse}Faire d’abord une étude commune puis demander une esquisse individuelle. La seconde dérivée sert ici au critère d’inflexion; ne pas exiger un cours général de convexité.\end{notepourlaclasse}
